\section{Auswertung des Jahresabschlusses}
\label{sec:abschluss-azn}

Auch dieser Abschnitt berichtet nur. Jede Kennzahl nennt ihre Gleichung und
rechnet sie vor; ein Ergebnis ohne seinen Rechenweg ist nicht pr\"ufbar.
Alle Betr\"age in Mio.\,USD, Gesch\"aftsjahr gleich Kalenderjahr.

\subsection{Umsatz}

Der Gesamtumsatz stieg \Pct{\AznUmsatzWachstum} auf \MioUSD{\AznUmsatzXXV}
nach \MioUSD{\AznUmsatzXXIV} im Vorjahr.\quelleGB{2025}{125}\quelleMerken{s8gbxcacfxbcf}
Er setzt sich zusammen aus Produktverk\"aufen von \MioUSD{\AznProduktumsatz},
Allianzertr\"agen von \MioUSD{\AznAllianzumsatz} und
Kooperationsertr\"agen.\quelleErneut{s8gbxcacfxbcf}
\[
  \text{Umsatzwachstum} = \frac{\AznUmsatzXXV}{\AznUmsatzXXIV} - 1
  = \Pct{\AznUmsatzWachstum}
\]

Auf die Vereinigten Staaten entfallen \MioUSD{\AznUmsatzUsa} und damit
\Pct{\AznUsaAnteil} des Umsatzes, auf Europa \MioUSD{\AznUmsatzEuropa}, auf die
Schwellenl\"ander \MioUSD{\AznUmsatzSchwellen}.\quelleGB{2025}{32}\quelleMerken{s8gbxcacfxdc}
China allein steuert \MioUSD{\AznUmsatzChina} bei, \Pct{\AznChinaAnteil} des
Umsatzes.\quelleGB{2025}{33}\quelleMerken{s8gbxcacfxdd}

\subsection{Der Fr\"uhindikator: Zulassungen statt Auftragseingang}

Ein Pharmaunternehmen f\"uhrt kein Auftragsbuch. An die Stelle des
Auftragseingangs tritt die Zahl der Projekte, die die Sp\"atphase erreichen,
denn sie bestimmt den Umsatz der Jahre nach dem Patentablauf. Der Bericht
z\"ahlt \Stueck{\AznPipeline} Projekte in der Entwicklung, davon
\Stueck{\AznSpaetphase} in Zulassungsstudien oder in der
Zulassungspr\"ufung.\quelleGB{2025}{30}\quelleMerken{s9gbxcacfxda}
Im Berichtsjahr gab es \Stueck{\AznZulassungen} Zulassungen neuer Wirkstoffe
und wesentlicher Indikationserweiterungen.\quelleGB{2025}{28}\quelleMerken{s9gbxcacfxci}
Der Forschungsaufwand betrug \MioUSD{\AznForschung} und damit
\Pct{\AznForschungsquote} des Umsatzes; bereinigt weist das Unternehmen
\MioUSD{\AznForschungCore} aus.\quelleGB{2025}{55}\quelleMerken{s9gbxcacfxff}

\subsection{EBITDA-Marge, ausgewiesen und bereinigt}

Das Unternehmen weist das EBITDA selbst aus: \MioUSD{\AznEbitda} nach
\MioUSD{\AznEbitdaVJ} im Vorjahr.\quelleGB{2025}{59}\quelleMerken{s9gbxcacfxfj}
Bezogen auf den Umsatz sind das \Pct{\AznEbitdaMarge} nach
\Pct{\AznEbitdaMargeVJ}.
\[
  \text{EBITDA-Marge} = \frac{\AznEbitda}{\AznUmsatzXXV} = \Pct{\AznEbitdaMarge}
\]

Daneben steht die bereinigte Gr\"o{\ss}e, die das Unternehmen \enquote{Core}
nennt: ein Betriebsergebnis von \MioUSD{\AznCoreBetriebsergebnis} gegen
\MioUSD{\AznBetriebsergebnisXXV} ausgewiesen.\quelleGB{2025}{51}\quelleMerken{s9gbxcacfxfb}
Die Differenz von rund einem Drittel entsteht aus Abschreibungen auf erworbene
Rechte, Restrukturierungsaufwand und Bewertungs\"anderungen bedingter
Kaufpreise. Das Verh\"altnis kehrt sich in keinem der beiden Jahre um: die
bereinigte Gr\"o{\ss}e liegt stets \"uber der ausgewiesenen. Getragen wird in
diesem Bericht die \emph{ausgewiesene} Gr\"o{\ss}e, weil sie den Zahlungen
folgt; die bereinigte steht daneben, weil die Prognose des Unternehmens in ihr
formuliert ist.

\subsection{Operatives Ergebnis}

Das ausgewiesene Betriebsergebnis betr\"agt \MioUSD{\AznBetriebsergebnisXXV},
eine Marge von \Pct{\AznBetriebsMarge}.\quelleGB{2025}{125}\quelleMerken{s9gbxcacfxbcf}
Der Restrukturierungsaufwand des Jahres lag bei
\MioUSD{\AznRestrukturierung}.\quelleErneut{s9gbxcacfxfj}

\subsection{Nettogewinn}

Der Jahres\"uberschuss stieg \Pct{\AznErgebnisWachstum} auf
\MioUSD{\AznJahresueberschussXXV}.\quelleErneut{s9gbxcacfxbcf}
Je Aktie sind das \USDje{\AznErgebnisJeAktieXXV} ausgewiesen und
\USDje{\AznCoreErgebnisJeAktie} bereinigt.\quelleErneut{s9gbxcacfxfb}
Die Eigenkapitalrendite betr\"agt
\[
  \frac{\AznJahresueberschussXXV}{\AznEigenkapital} = \Pct{\AznEigenkapitalrendite},
\]
gerechnet auf das Eigenkapital zum Bilanzstichtag von
\MioUSD{\AznEigenkapital}.\quelleGB{2025}{126}\quelleMerken{s9gbxcacfxbcg}
Die Eigenkapitalquote liegt bei \Pct{\AznEigenkapitalquote}, die
Nettoverschuldung bei \MioUSD{\AznNettoverschuldung} und damit beim
\Faktor{\AznVerschuldungsgrad}-fachen des EBITDA.\quelleGB{2025}{60}\quelleMerken{s9gbxcacfxga}
Net Debt zu EBITDA setzt die Nettoverschuldung ins Verh\"altnis zum
Betriebsergebnis \emph{vor} Abschreibungen, nicht zum Gewinn.

\subsection{Dividende und Dividendenrendite}

F\"ur das Jahr 2025 wurden \USDje{\AznDividendeJeAktie} je Aktie erkl\"art,
gezahlt in zwei Raten.\quelleGB{2025}{63}\quelleMerken{s10gbxcacfxgd}
Am Bilanzstichtag ausgezahlt und verbucht wurden
\USDje{\AznDividendeGezahltJeAktie}. Die Differenz ist kein Widerspruch,
sondern der Versatz zwischen Erkl\"arung und Zahlung: die zweite Rate eines
Jahres flie{\ss}t erst im M\"arz des folgenden.\quelleGB{2025}{169}\quelleMerken{s10gbxcacfxbgj}
Bei einem Kurs von \USDje{\AznKurs} betr\"agt die Dividendenrendite
\Pct{\AznDividendenrendite}.

\subsection{Aktienkursverlauf}
\label{sec:kursbild-azn}

\begin{figure}[htbp]
\centering
\begin{tikzpicture}
\begin{axis}[kursachseazn]
\addplot+[gray!55, thick, mark=none] table[col sep=comma, x=datum,
  y=kurs] {data/kursverlauf_azn.csv};
\addplot+[kurspunkte] table[col sep=comma, x=datum,
  y=kurs, meta=art] {data/kursverlauf_azn.csv};
\end{axis}
\end{tikzpicture}
\caption{Kursverlauf der AstraZeneca-Aktie vom September 2016 bis zum
\Datum{\AznKursTag}. Jeder Punkt steht an seinem tats\"achlichen Handelstag:
{\color{kurshoch}$\blacktriangle$} Jahreshoch,
{\color{kurstief}$\blacktriangledown$} Jahrestief,
{\color{kursschluss}$\bullet$} Jahresschluss. Die Reihe ist auf die
Stammaktie zur\"uckgerechnet (siehe Abschnitt~\ref{sec:kurs-azn}).}
\label{fig:kurs-azn}
\end{figure}

\begin{table}[htbp]
\centering
\caption{Jahresschluss-, H\"ochst- und Tiefstkurse in USD. Das Jahr 2026 reicht
bis zum \Datum{\AznKursTag}. Sekund\"arquelle: Nasdaq.}
\label{tab:kurs-azn}
\begin{tabular}{lrrr}
\toprule
Jahr & Schluss & H\"ochst & Tiefst\\
\midrule
\tabellenkoerper{data/kurs_azn_tabelle}
\bottomrule
\end{tabular}
\end{table}

\FloatBarrier

\subsection*{Kennzahlen\"uberblick}
\addcontentsline{toc}{subsection}{Kennzahlen\"uberblick}

\begin{table}[htbp]
\centering
\caption{Kennzahlen 2025 gegen 2024, Betr\"age in Mio.\,USD}
\label{tab:azn-kennzahlen}
\begin{tabular}{lrrr}
\toprule
Kennzahl & 2025 & 2024 & Ver\"anderung\\
\midrule
Umsatz                       & \Betrag{\AznUmsatzXXV} & \Betrag{\AznUmsatzXXIV} & \Pct{\AznUmsatzWachstum}\\
Betriebsergebnis             & \Betrag{\AznBetriebsergebnisXXV} & \Betrag{\AznBetriebsergebnisXXIV} & \Pct{\AznBetriebsergebnisWachstum}\\
Betriebsergebnis (bereinigt) & \Betrag{\AznCoreBetriebsergebnis} & \Betrag{\AznCoreBetriebsergebnisVJ} & \Pct{\AznCoreWachstum}\\
EBITDA                       & \Betrag{\AznEbitda} & \Betrag{\AznEbitdaVJ} & \Pct{\AznEbitdaWachstum}\\
Jahres\"uberschuss           & \Betrag{\AznJahresueberschussXXV} & \Betrag{\AznJahresueberschussXXIV} & \Pct{\AznErgebnisWachstum}\\
Ergebnis je Aktie (USD)      & \Kurswert{\AznErgebnisJeAktieXXV} & \Kurswert{\AznErgebnisJeAktieXXIV} & \Pct{\AznEpsWachstum}\\
Freier Mittelzufluss         & \Betrag{\AznFcfXXV} & \Betrag{\AznFcfXXIV} & \Pct{\AznFcfWachstum}\\
Nettoverschuldung            & \Betrag{\AznNettoverschuldung} & \Betrag{\AznNettoverschuldungVJ} & \Pct{\AznNettoverschuldungWachstum}\\
Eigenkapital                 & \Betrag{\AznEigenkapital} & \Betrag{\AznEigenkapitalVJ} & \Pct{\AznEigenkapitalWachstum}\\
\bottomrule
\end{tabular}
\end{table}

\subsection{Die Mehrjahresreihe und die Lage im Zyklus}
\label{sec:reihe-azn}

Zwei Jahre sagen, was sich ver\"andert hat, und nichts dar\"uber, ob der
j\"ungste Wert normal ist. Genau daran h\"angt aber das Votum. Tabelle
\ref{tab:azn-reihe} f\"uhrt deshalb sieben Jahre; sie beginnt 2019, weil die
Gesch\"aftsberichte 2019 und 2020 in einer Schrift gesetzt sind, aus der sich
die gedruckte Seitenzahl nicht auslesen l\"asst, und ein Wert ohne
belegbare Seite in dieses Dokument nicht aufgenommen wird.

\begin{table}[htbp]
\centering
\caption{Sieben Jahre, Betr\"age in Mio.\,USD; freier Mittelzufluss je Aktie
in USD}
\label{tab:azn-reihe}
\begin{tabular}{lrrrrr}
\toprule
Jahr & Umsatz & Jahres\-\"uber\-schuss & operativer
     & freier & je Aktie\\
     &        &                        & Mittelzufluss & Mittelzufluss & \\
\midrule
\tabellenkoerper{data/flow_azn_reihevoll}
\bottomrule
\end{tabular}
\end{table}

Der freie Mittelzufluss ist hier definiert als operativer Mittelzufluss
abz\"uglich der Investitionen in Sachanlagen und immaterielle
Verm\"ogenswerte, zuz\"uglich der Erl\"ose aus deren Abgang und abz\"uglich der
Leasingtilgung. Die Leasingtilgung geh\"ort dazu, weil sie wirtschaftlich
Miete ist und dem Eigent\"umer nicht zur Verf\"ugung steht.

Zwei Ablesungen aus der Reihe, die der Bericht selbst nicht vornimmt. Erstens
wuchs der freie Mittelzufluss je Aktie von \USDje{\AznFcfJeAktieXIX} im Jahr
2019 auf \USDje{\AznFcfVoll} im Jahr 2025; das sind \Pct{\AznCagr} pro Jahr
\"uber sechs Jahre. Zweitens ist \USDje{\AznFcfVoll} zugleich der \emph{h\"ochste}
Wert der Reihe; das Jahr 2025 liegt im hundertsten Perzentil seiner eigenen
sieben Jahre. Der Median betr\"agt \USDje{\AznFcfMedian}, die letzten zw\"olf
Monate bis zum 30.~Juni 2026 ergeben \USDje{\AznFcfLtm}.
Abschnitt~\ref{sec:flow-azn} rechnet mit allen dreien und sagt bei jedem
Ergebnis dazu, welche Zeile es tr\"agt.
